%% =====================================================================
%%  要旨
%%  研究の背景・課題・目的・方法・結果・結論を 1 ページ（800〜1200 字程度）にまとめる．
%%  要旨だけを読んでも研究の内容がわかるように書く．
%%  図・表・数式・文献の引用は使わない．
%%  下の例の「%% 背景」などは構成の目安を示すコメント（PDF には出ない）．
%%  「（背景）」のような見出しを本文に書く必要はない．
%%  修士論文の場合，研究科に別途提出する「論文要旨」（2〜8 頁，Hoppii の様式）は，
%%  この要旨とは別の書類である．
%% =====================================================================
\begin{abstractJA}
\begin{writingnote}
要旨は本文を書き終えてから最後に書き，本文にない内容は入れない．
欠かせないのは目的と主な結果である．背景と課題は 1〜2 文にとどめ，方法と結果に字数を割く．
結果は「有効性を確認した」のような抽象的な表現で済ませず，どのような条件で何がわかったかを具体的に書く（主要な数値を示すとよい）．
最後に，結果から言えることと，提案手法の適用条件や限界を述べる．
略語は初出で正式な名称を示す．用語と数値は本文，特に「結論」の章と一致させる．
\end{writingnote}

%% 背景
近年，……の分野では，○○が重要になっている．
%% 課題
しかし，従来の方法には……という課題がある．

%% 目的・方法
そこで本研究では，……を目的として，○○を用いた△△問題の解法を提案する．
提案手法では，……により……する．

%% 結果・結論
数値実験により，……の条件で……という結果が得られた．
以上より，提案手法の適用条件として……，限界として……が明らかになった．
\end{abstractJA}

%% 英文の要旨（Abstract）も付ける場合は，次の行頭の % を外す．
% \begin{abstractEN}
% Write the abstract in English here.
% \end{abstractEN}
